% =================================================================
%  style/cover.tex
%  Обложка, страница составителей и оглавление
%
%  Использование в теле документа:
%     \makecover{AnalysisColor}{Математический анализ}%
%               {Билеты}{Фамилия И.\,О.}
%     \makestaffandcontents
% =================================================================

\usepackage{tikz}
\usetikzlibrary{calc}
\usepackage{xcolor}

% -----------------------------------------------------------------
% Цвета обложек по предметам
% -----------------------------------------------------------------
\definecolor{AlgosColor}{HTML}{15162B}     % алгоритмы
\definecolor{AlgebraColor}{HTML}{301010}   % алгебра
\definecolor{AnalysisColor}{HTML}{102619}  % анализ
\definecolor{DiscreteColor}{HTML}{231030}  % дискретная математика
\definecolor{MetaColor}{HTML}{6672B8}
\definecolor{LineColor}{HTML}{D8D8E6}
\definecolor{SoftText}{HTML}{7A7A92}
\definecolor{SoftGray}{HTML}{CFCFE6}

% -----------------------------------------------------------------
% Шапка обложки и город/год — правьте под свой курс
% -----------------------------------------------------------------
\newcommand{\coverheader}{СПбГУ \ $\cdot$ \ МКН \ $\cdot$ \ СП \ $\cdot$ \ 1 КУРС}
\newcommand{\coverplaceyear}{Санкт-Петербург, 2026}

% -----------------------------------------------------------------
% \makecover{<цвет фона>}{<название курса>}{<подзаголовок>}{<лектор>}
% -----------------------------------------------------------------
\newcommand{\makecover}[4]{%
	\thispagestyle{empty}
	
	\begin{titlepage}
		\begingroup
		\sffamily
		
		\begin{tikzpicture}[remember picture, overlay]
			\fill[#1] (current page.south west) rectangle (current page.north east);
			
			% Диагональная «сетка» поверх фона
			\begin{scope}[white, line width=0.4pt]
				\foreach \i in {0,1,...,25} {
					\draw[opacity=0.07] ($(current page.north west)+(0,-\i*1.2)$) --
					($(current page.north east)+(-\i*2,-28)$);
				}
				\foreach \i in {0,1,...,20} {
					\draw[opacity=0.05] ($(current page.south east)+(0,\i*1.8)$) --
					($(current page.south west)+(\i*3.5,18)$);
				}
				\foreach \i in {0,5,...,30} {
					\draw[opacity=0.15] ($(current page.north east)+(-\i,0)$) --
					($(current page.south west)+(0,\i)$);
				}
			\end{scope}
			
			% Верхняя строка
			\node[anchor=north west, white, inner sep=0pt]
			at ($(current page.north west)+(2cm,-1.5cm)$) {%
				\fontsize{16}{18}\selectfont\bfseries \coverheader
			};
			
			% Заголовок и подзаголовок
			\node[anchor=west, white, inner sep=0pt, text width=18cm, align=left]
			at ($(current page.west)+(2cm,5cm)$) {%
				{\fontsize{52}{58}\fontseries{bx}\selectfont #2\par}
				\vspace{0.7cm}
				{\fontsize{30}{36}\selectfont\bfseries #3\par}
			};
			
			% Лектор
			\node[anchor=west, white, inner sep=0pt, text width=10cm, align=left]
			at ($(current page.west)+(2cm,1.1cm)$) {%
				{\fontsize{20}{20}\selectfont\bfseries Лектор}
				{\fontsize{22}{24}\selectfont #4}
			};
			
			% Город и год
			\node[anchor=south, white, inner sep=0pt]
			at ($(current page.south)+(0,2cm)$) {%
				{\fontsize{18}{20}\selectfont \coverplaceyear}
			};
		\end{tikzpicture}
		\endgroup
	\end{titlepage}
	
	\clearpage
}

% -----------------------------------------------------------------
% Карточки составителей
%   \leadperson{<имя>}{<github>}{<роль>}   — крупная карточка с ролью
%   \nameghblock{<имя>}{<github>}          — имя + ссылка на github
%   \nameonlyblock{<имя>}                  — только имя
% -----------------------------------------------------------------
\newcommand{\leadperson}[3]{%
	\begin{minipage}[t]{0.42\textwidth}
		\centering
		\parbox[t][2.7em][t]{0.95\linewidth}{%
			\centering
			{\fontsize{10.5}{12.5}\selectfont\color{SoftText} #3\par}
		}%
		
		\vspace{0.08cm}
		{\fontsize{15.5}{18}\selectfont #1\par}
		\vspace{0.10cm}
		{\fontsize{12}{14}\selectfont\color{MetaColor}\href{https://github.com/#2}{@#2}\par}
	\end{minipage}%
}

\newcommand{\nameghblock}[2]{%
	\begin{minipage}[t]{0.30\textwidth}
		\centering
		{\fontsize{14.5}{17}\selectfont #1\par}
		\vspace{0.08cm}
		{\fontsize{11.5}{13}\selectfont\color{MetaColor}\href{https://github.com/#2}{@#2}\par}
	\end{minipage}%
}

\newcommand{\nameonlyblock}[1]{%
	\begin{minipage}[t]{0.30\textwidth}
		\centering
		{\fontsize{14.5}{17}\selectfont #1\par}
	\end{minipage}%
}

% -----------------------------------------------------------------
% Страница «Составители» + оглавление.
% Список людей правится прямо здесь.
% -----------------------------------------------------------------
\newcommand{\makestaffandcontents}{%
	\clearpage
	\thispagestyle{empty}
	\begingroup
	\sffamily
	
	{\centering
		{\fontsize{22}{26}\selectfont\bfseries Составители\par}
		\par}
	
	\vspace{0.72cm}
	
	% --- ведущие роли ---
	\begin{center}
		\leadperson{Имя Фамилия}{github-login}{координация проекта, техническая редакция, вёрстка}
		\hspace{0.05\textwidth}
		\leadperson{Имя Фамилия}{github-login}{проверка и научное редактирование}
	\end{center}
	
	\vspace{0.62cm}
	
	% --- остальные участники: по одному, по два или по три в ряд ---
	\begin{center}
		\nameghblock{Имя Фамилия}{github-login}
	\end{center}
	
	\vspace{0.92cm}
	
	% --- разделитель ---
	\noindent{\color{LineColor}\rule{\textwidth}{0.55pt}}
	
	\vspace{0.62cm}
	
	{\fontsize{13}{16}\selectfont
		\setcounter{tocdepth}{1}
		\tableofcontents
	}
	
	\vfill
	
	{\centering
		{\fontsize{11.5}{13}\selectfont\color{SoftText} \coverplaceyear\par}
		\par}
	
	\endgroup
	\clearpage
}
